\chapter{LITERATURE REVIEW}

The development of AI technology in talent acquisition has been extensively researched and discussed both in academic research and whitepapers by the industry. Current research in the period from 2024 to 2026 shows that there has been a clear move beyond mere text parsing to semantic understanding, reducing biases, and algorithmic explainability. This chapter provides a review of current state-of-the-art technologies and practices in several different aspects of operation to frame the prototypical solution to be presented herein.

\section{Traditional Keyword-Based Systems}
Early approaches to resume parsing and applicant tracking system (ATS) software relied on rule-based algorithms and Boolean text searches. Recent comparisons of industry practices have shown that modern ATS systems perform extraction of text and calculate the occurrence of keywords. Though being computationally inexpensive and easy to implement, such an approach is not effective as it results in very high numbers of false negatives. The use of synonyms, different formulations, and descriptive language ("orchestrated a financial turnaround" as opposed to "leadership") prevents many candidates from passing even though they meet the requirements for the position. Furthermore, modern systems are not robust enough against manipulations; applicants can simply fill the text with keywords or hidden text to achieve higher scores.

\section{Machine Learning and TF-IDF Models}
In order to go beyond Boolean logic while avoiding deep learning's heavy computing load and associated costs, it became common practice to merge classical ML techniques with statistical NLP approaches. Libraries like SpaCy used for Named Entity Recognition (NER) in combination with TF-IDF showed promising results in performing initial feature extraction and structured parsing tasks. The latter is defined as the measure of significance for each term in relation to the whole body of text and penalizes very popular terms, making it harder to do keyword stuffing naively. Combining TF-IDF with distance measures, including Cosine Similarity, enables calculating relevance score between job description and résumé vectors, providing much higher precision compared to plain keyword matching, yet still running in milliseconds on a standard CPU.

\section{Transformer and Semantic Matching Models}
Perhaps the greatest technical development that came in automated résumé screening over the last three years is the use of transformers. Empirical studies comparing BERT models with Sentence-BERT prove its unprecedented ability to put human language into context. Deshmukh and Raut (2025) found that NLP models built with BERT can significantly boost screening speed by understanding that "Python" and "Django" belong to the same semantic field, hence solving the keyword mismatch problem. Along with that, the Resume2Vec system (2025) introduced an innovative way to screen candidates based on their skills represented as multi-dimensional vector in space, showing remarkable improvements in terms of Normalized Discounted Cumulative Gain (nDCG) metric compared to the standard ATSs. Finally, recent works on multi-agent LLMs, such as Lo et al. (2025)'s system using RAG method, showed competitive precision in résumé screening according to different requirements, being nearly equal in accuracy to HR professionals.

\section{Explainable AI (XAI) in Recruitment}
Inasmuch as artificial intelligence takes up an increasing role in employment regulation, complexity of the algorithm used is a notable concern that poses significant risks. The "black-box" nature of neural networks is a drawback for recruiters in terms of being able to determine why some candidates receive a 45\% while others an 85\% matching score. Explainable AI (XAI) attempts to address this problem by making the model more comprehensible. Much effort has been made toward developing post-hoc explanations, including such methodologies as SHAP (Shapley Additive Explanations) and LIME (Local Interpretable Model-agnostic Explanations). Mathematical methods try to reconstruct the model and find out which specific criteria were used to classify candidates (i.e., years of experience compared to certain skills). Despite their accuracy, these methods are computationally costly and require extensive knowledge to be interpreted properly. A gap in the existing literature can be filled with studies into rule-based XAI output which produces clear explanations in simple English specifically for HR specialists.


\section{Bias, Fairness, and Format Discrimination}
The problem of algorithmic bias is another major topic in recruiting automation research. While many works have been done in relation to issues such as demographic discrimination and various biases related to factors like gender, race, or age, there is another problem that needs more attention. It is known as "Format Bias" or "Layout Bias," and automated parsers usually suffer from catastrophic errors in processing unusual formats that contain many columns and lots of visuals. In other words, when job seekers try to impress the employer using sophisticated formatting to make themselves stand out, they may face additional drawbacks due to limitations associated with automatic processing as a result of the failure to recognize boundaries of text streams. Recent advancements in layout analysis prove that the development of robust OCR fallbacks (Tesseract, for example), as well as spacing hierarchy recognition technology, can be crucial in this regard.


\section{Cost, Scalability, and Deployment Constraints}
While large language models and Sentence-BERT models show impressive academic performance concerning semantic discrepancies between terms used by recruiters and job seekers, their practical application is severely restricted by cost and computational power issues. To process tens of thousands of resumes at once, models containing 7 billion or 70 billion parameters require powerful memory resources and expensive cloud GPU instances, which makes top-tier AI solutions out of reach for smaller organizations. As a consequence, another market niche is rapidly developing, requiring lightweight, inexpensive solutions that would not involve expensive hardware. While such systems are less advanced regarding semantic interpretation, they work effectively as primary screening filters.

\section{Literature Comparison}
\begin{table}[H]
\centering
\caption{Comparison of Recent Literature in AI Resume Screening}
\begin{tabularx}{\textwidth}{p{0.18\textwidth}Y Y Y Y}
\toprule
\textbf{Author / Year} & \textbf{Method and Technology} & \textbf{Focus Area / Strengths} & \textbf{Limitations} & \textbf{Relevance to Prototype}\\
\midrule
Deshmukh and Raut (2025) & BERT-based NLP vectorization & High semantic matching accuracy and precision mapping. & Computationally heavy; acts as a black box without external XAI. & Highlights the benchmark semantic gap our system evaluates against.\\
Lo et al. (2025) & Multi-Agent RAG-LLM & Context-aware, adapts to diverse hiring roles dynamically. & High API costs, excessive latency for bulk resume screening workflows. & Demonstrates the ceiling of semantic AI; validates our cost-constraint argument.\\
Jamal et al. (2024) & SHAP and LIME integration & Unravels black-box decisions, increasing HR trust in algorithms. & Complex mathematical outputs are not easily read by laymen recruiters. & Supports the necessity of our plain-English natural language XAI module.\\
Jaha (2024) & Semantic Document Layout Analysis & Extracts text accurately from complex handwritten or scanned structures. & High computational overhead for preprocessing standard digital PDFs. & Reinforces the need for our Tesseract OCR fallback pipeline logic.\\
Alsubaie and Aleisa (2025) & XAI Bias Mitigation via Random Forest and XGBoost & Identifies and isolates biased features, such as demographics, systematically. & Focuses strictly on structured data rather than raw unstructured resume text. & Highlights the ethical mandate for transparency and fairness in scoring models.\\
\bottomrule
\end{tabularx}
\end{table}

\section{Identified Research Gaps}
According to the current body of literature, there is a clear divergence between technologies at two extremes. On the one side, conventional keyword-based matchers provide fast execution and minimal costs at the expense of accuracy and vulnerability to manipulation. On the other hand, sophisticated large language models achieve great levels of accuracy; however, they come with significant costs, considerable latency, and lack transparency. Thus, the key research challenge lies within the middle zone, which involves devising an approach that is systematic, interpretable, and computationally economical. Ideally, the method needs to go beyond the elementary task of identifying keywords by means of TF-IDF and NLP approaches, include proper mathematical measures to prevent manipulation attacks, process documents using optical character recognition to avoid layout bias, and generate transparent “white box” interpretations in plain English language in milliseconds on a regular CPU.
